\section*{Part III: Parcel Diligence and How to Buy Land}
\addcontentsline{toc}{section}{Part III: Parcel Diligence and How to Buy Land}


\subsection{16. A letter of intent is not title}

Do not begin with a large non-refundable token payment. Obtain identifiers and records; appoint independent legal counsel with local conveyancing expertise; establish the seller's authority; investigate title and encumbrances; verify use; and only then sign a carefully conditional agreement with clear refund and exit rights. Transaction sequencing matters because sellers may seek early commitment before the buyer understands restrictions.

The legal ability to buy a plot is distinct from the legal ability to use it for a university, research laboratory, hostel, or campus.


\subsection{17. Parcel identification in Maharashtra}

Use the survey/gat number, city survey or CTS number, boundaries, village, taluka, district, and georeferenced map. Maharashtra's Land Records Department provides online 7/12 extracts, 8A, property cards, mutation records, maps and archival services [P1]. These records are essential starting points but \textbf{are not a self-sufficient certificate of unencumbered legal title}.

Match records to physical possession through a licensed survey and boundary inspection. Ask for historical ownership deeds, partitions, inheritance documents, court proceedings, mortgages and release deeds, and acquisition or reservation notices as relevant.


\subsection{18. Title search and encumbrances}

Have counsel review an appropriate historical chain of title, registered instruments, Index II or other registration records, mortgages/charges, easements, liens, restrictions, litigations, family claims, tenancy and possession, government dues, and encroachments. Search court and revenue proceedings where relevant. The review period and required documents should be advised by local counsel; a superficial ``clear 7/12'' is not enough.

A title report should identify unresolved risks, the proposed contractual cure, and whether title is sufficiently marketable for the intended lender and use.


\subsection{19. Agricultural land, conversion and restricted transfers}

In Maharashtra, land classification, agricultural status, permitted transferees, tribal land restrictions, ceiling provisions, conversion and planning rules can affect whether a particular buyer may acquire or develop a parcel. These are fact-intensive. A registered company should not assume it can buy rural agricultural land for a campus simply because a broker says ``NA conversion is easy.''

Get a parcel-specific legal opinion identifying permissions needed \textbf{before acquisition}, before construction, and before educational occupation. Do not budget on speculative approvals as if they are secured.


\subsection{20. Planning and zoning}

Obtain current regional/development plan extracts, local planning authority jurisdiction, land-use reservations, proposed roads, permissible uses, floor space index or FAR where applicable, setbacks, height limits, parking, building-safety requirements, and any special control regulations. Check environmental and infrastructure constraints. Map land price into \textbf{legally buildable and occupiable academic space}.

Approval of a building as ``commercial'' or ``institutional'' in 1 sense does not automatically make it suitable for a recognized degree institution or particular research activity.


\subsection{21. Access and easements}

Confirm legally enforceable road access and its width, rights over private approach roads, rights of way, construction vehicle access, emergency-vehicle access, and availability of fire tender turning areas. A plot reachable by informal permission over a neighbour's land may be unusable for financing and approvals.

Verify all access on site. If an access easement must be granted or registered, make it a condition precedent to closing, not a verbal reassurance.


\subsection{22. Water, power, drainage and telecom}

Research facilities may require substantial power quality, backup generation, cooling, water, treatment, drainage, waste segregation, and high-availability network links. A seller's claim that a transformer is nearby is not a capacity commitment from a utility. Obtain written feasibility or estimates from relevant providers, quantify connection lead times, and include off-site infrastructure in the cost plan.

For a potential residential campus, water security, wastewater treatment and public-health systems can dominate operating risk. Local environmental constraints may be decisive regardless of land affordability.


\subsection{23. Geotechnical, flood and environmental diligence}

Commission topographic survey, soil/geotechnical investigation, flood-risk review, drainage, slope stability and relevant environmental screening before fixing the building budget. A cheap plot with weak soil or poor drainage can require expensive foundations and land preparation. Wetland, forest, protected-area or pollution-related restrictions may render proposed development unlawful or costly.

If there are prior industrial uses, investigate contamination. Require specialist reports and allocate remediation obligations contractually where necessary.


\subsection{24. Land assembly and holdouts}

A campus may need adjacent parcels. Each additional owner creates price, title, timing and holdout risk. Do not assume a master plan can be implemented when 1 critical plot remains unsigned. Consider right-sized initial phases with legally independent access and services. In some situations a smaller contiguous plot with clean title is superior to a theoretically larger but fragmented assembly.

Secure options or conditional agreements only with proper legal protections and realistic expiry periods; the market can move while approvals are pending.


\subsection{25. Purchase agreement: essential architecture}

A serious agreement should address exact property description and annexed survey, price and payment milestones, seller warranties, representations as to title and litigation, encumbrance release, possession, access and utilities, conditions precedent, required regulatory consents, due diligence period, indemnities, escrow mechanics, tax and statutory charges, default remedies, long-stop date and dispute resolution. Obtain advice on stamp duty and registration; obligations may arise before completion.

A founder should never substitute a broker's 1-page form for independently reviewed conveyancing documentation on a major institutional asset.


\subsection{26. Escrow, finance and closing controls}

Prefer documented movement of consideration with appropriate escrow and lien-release procedures. Confirm the bank's conditions for mortgageability, title defects, charge registration, and equity contribution before signing hard commitments. Verify seller KYC, authority to execute, board approval for company sellers, and any prior lender's no-objection or release requirements.

At closing reconcile payments, registered instruments, possession memorandum, original document custody, mutation applications, taxes, and insurance transition. Physical occupation without legally completed title can become a decades-long dispute.


\subsection{27. Post-acquisition actions}

After registration, complete mutation and relevant property-card/land-record updates, secure possession, install lawful boundaries and security, verify taxes and insurance, protect the site from encroachment and dumping, and track approval conditions. Mutation is an administrative record update, not a substitute for the original conveyance and proper title opinion.

Maintain a digital and physical property dossier accessible to authorized management and auditors, not only to the founder or an external broker.


\subsection{28. Acquisition red flags}

Hard-stop or major-review signals include unclear title chains, unregistered transfers, disputed inheritance, unresolved litigation, encroachments, illegal road access, prohibited transfers, unexpected government reservations, inability to use for institutional purposes, adverse environmental findings, non-cooperative sellers, verbal-only utility assurances, pressure for cash payments, and financing contingent on assumptions a lender has not accepted.

Red flags should be resolved by evidence, contract, or withdrawal-not by applying a smaller arbitrary risk discount to an otherwise attractive score.
